\section{本讲小结}

仪表是电路中的元件，量程改装应从表头满偏电流和内阻出发。测量改变了被测电路，读数属于接入仪表后的真实节点、支路。伏安法应分别说明分流和分压，并在需要时用已知内阻修正。

欧姆表的电阻刻度由内部电源和非线性电流关系产生。倍率、欧姆调零、机械零点及表笔极性各有含义，不能互相替代。故障判断应沿实际导电路径逐段核查。非线性元件需在各工作点计算电阻、功率，实验曲线和拟合还应保留读数精度及模型条件。

